\documentclass[11pt,a4paper]{article}

\usepackage[utf8]{inputenc}
\usepackage[T1]{fontenc}
\usepackage[french]{babel}
\usepackage{lmodern}
\usepackage{geometry}
\geometry{margin=2.2cm}
\usepackage{graphicx}
\usepackage{booktabs}
\usepackage{array}
\usepackage{hyperref}
\usepackage{xcolor}
\usepackage{amsmath}
\usepackage{float}
\usepackage{listings}

\hypersetup{
  colorlinks=true,
  linkcolor=blue,
  urlcolor=blue,
  citecolor=blue
}

\title{\textbf{Pr\'ediction du churn client dans les t\'el\'ecommunications}\\\vspace{0.3cm}\large Rapport de projet Machine Learning (End-to-End)}
\author{Otmane ZIRARI}
\date{F\'evrier 2026}

\begin{document}
\maketitle

\begin{abstract}
Ce rapport pr\'esente un projet complet (end-to-end) de pr\'ediction du \emph{churn} client dans le secteur des t\'el\'ecommunications \`a partir du jeu de donn\'ees \emph{Telco Customer Churn}. L'objectif est de d\'evelopper un pipeline reproductible incluant le pr\'etraitement des donn\'ees, l'entra\^inement de plusieurs mod\`eles, l'\'evaluation via des m\'etriques pertinentes pour un probl\`eme d\'es\'equilibr\'e, ainsi qu'un d\'eploiement d'une interface web permettant l'inf\'erence pour de nouveaux clients.
\end{abstract}

\tableofcontents
\newpage

\section{Introduction et contexte}
Le \textbf{churn client} correspond au fait qu'un client quitte un service (r\'esiliation, non-renouvellement). Dans les t\'el\'ecommunications, la concurrence est forte et le co\^ut d'acquisition d'un nouveau client est g\'en\'eralement sup\'erieur au co\^ut de r\'etention. Il est donc strat\'egique d'identifier en amont les clients \`a risque afin de d\'eclencher des actions de r\'etention (offres cibl\'ees, support, engagement).

Ce projet s'inscrit dans une d\'emarche de \textbf{classification binaire}: pr\'edire si un client va \textbf{churn} (classe 1) ou \textbf{ne pas churn} (classe 0), \`a partir de variables d\'emographiques, contractuelles et de consommation.

\section{Objectifs du projet}
\begin{itemize}
  \item Concevoir un pipeline complet de Machine Learning: donn\'ees $\rightarrow$ pr\'etraitement $\rightarrow$ mod\`eles $\rightarrow$ \'evaluation.
  \item Comparer plusieurs algorithmes (lin\'eaires, arbres, ensembles, boosting, SVM).
  \item G\'erer le d\'es\'equilibre des classes (churn minoritaire) via pond\'eration et SMOTE.
  \item Sauvegarder les r\'esultats (m\'etriques, matrice de confusion, rapport de classification) pour chaque exp\'erience.
  \item D\'eployer une interface web (Flask) permettant \`a un utilisateur de saisir les informations d'un client et obtenir une pr\'ediction.
\end{itemize}

\section{Jeu de donn\'ees}
\subsection{Source}
Le dataset utilis\'e est \emph{Telco Customer Churn} (Kaggle):
\begin{itemize}
  \item Taille: \textasciitilde 7043 clients.
  \item Cible: \texttt{Churn} (Yes/No).
  \item Variables: d\'emographie, services, contrat, facturation.
\end{itemize}

\subsection{Variables principales}
Exemples de variables:
\begin{itemize}
  \item \textbf{D\'emographie}: \texttt{gender}, \texttt{SeniorCitizen}, \texttt{Partner}, \texttt{Dependents}.
  \item \textbf{Services}: \texttt{PhoneService}, \texttt{InternetService}, \texttt{TechSupport}, \texttt{StreamingTV}...
  \item \textbf{Contrat/facturation}: \texttt{tenure}, \texttt{Contract}, \texttt{PaymentMethod}, \texttt{MonthlyCharges}, \texttt{TotalCharges}.
\end{itemize}

\subsection{D\'es\'equilibre de classes}
La classe \emph{churn} repr\'esente environ 27\% des clients. Ce d\'es\'equilibre peut induire des mod\`eles qui maximisent l'accuracy mais d\'etectent mal la classe minoritaire. Ainsi, des m\'etriques comme \textbf{Recall}, \textbf{F1-score} et \textbf{ROC-AUC} sont essentielles.

\section{Architecture du projet et organisation des fichiers}
\subsection{Structure}
\begin{verbatim}
archivedata/
  data/WA_Fn-UseC_-Telco-Customer-Churn.csv
  src/preprocessing.py
  src/train.py
  src/evaluate.py
  outputs/preprocessed/
  outputs/runs/
  models/best_model.pkl
  templates/index.html
  app.py
  requirements.txt
\end{verbatim}

\subsection{R\^ole des scripts}
\begin{itemize}
  \item \texttt{src/preprocessing.py}: pr\'etraitement (nettoyage, encodage, split, standardisation).
  \item \texttt{src/train.py}: entra\^inement multi-mod\`eles, sauvegarde des m\'etriques/artifacts, (optionnel) MLflow.
  \item \texttt{src/evaluate.py}: visualisations (ROC, PR), matrice de confusion, rapports.
  \item \texttt{app.py} + \texttt{templates/index.html}: interface web et API d'inf\'erence.
\end{itemize}

\section{M\'ethodologie}
\subsection{Pr\'etraitement des donn\'ees}
Les \`etapes principales impl\'ement\'ees dans \texttt{preprocessing.py} sont:
\begin{enumerate}
  \item \textbf{Chargement} du CSV.
  \item \textbf{Nettoyage} de \texttt{TotalCharges}: conversion num\'erique; valeurs invalides $\rightarrow$ 0.
  \item \textbf{Encodage}:
  \begin{itemize}
    \item Binaire Yes/No $\rightarrow$ 1/0 (ex: \texttt{Partner}, \texttt{Dependents}, \texttt{Churn}).
    \item Variables multi-cat\'egories encod\'ees via \texttt{LabelEncoder}.
  \end{itemize}
  \item \textbf{S\'eparation} features/target (exclusion \texttt{customerID} et \texttt{Churn}).
  \item \textbf{Split} train/test stratifi\'e.
  \item \textbf{Standardisation} optionnelle avec \texttt{StandardScaler}.
\end{enumerate}

\subsection{Mod\`eles entra\^in\'es}
Les mod\`eles entra\^in\'es dans \texttt{train.py}:
\begin{itemize}
  \item \textbf{Logistic Regression} (baseline, interpr\'etable).
  \item \textbf{Decision Tree} (r\`egles de d\'ecision, interpr\'etable).
  \item \textbf{Random Forest} (ensemble d'arbres, robuste).
  \item \textbf{XGBoost} (boosting, souvent performant).
  \item \textbf{LightGBM} (boosting, rapide et performant).
  \item \textbf{SVM} (fronti\`ere non-lin\'eaire, kernel RBF).
\end{itemize}

\subsection{Gestion du d\'es\'equilibre}
Deux strat\'egies ont \`et\'e compar\'ees:
\begin{itemize}
  \item \textbf{Pond\'eration des classes} (\texttt{class\_weight='balanced'}).
  \item \textbf{SMOTE}: sur-\'echantillonnage synth\'etique de la classe minoritaire sur le train set.
\end{itemize}

\subsection{M\'etriques d'\'evaluation}
Les m\'etriques calcul\'ees:\newline
\begin{itemize}
  \item Accuracy
  \item Precision
  \item Recall
  \item F1-score
  \item ROC-AUC
\end{itemize}
Dans un contexte de churn, le \textbf{Recall} est important (ne pas rater des churners), tandis que le \textbf{F1-score} donne un compromis.

\section{Suivi des exp\'eriences avec MLflow}
\subsection{Objectif}
MLflow est utilis\'e comme outil de \textbf{tracking d'exp\'eriences} afin de conserver une trace des entra\^inements, comparer les mod\`eles et assurer la reproductibilit\'e.

\subsection{Ce qui est enregistr\'e}
Lors de l'ex\'ecution de \texttt{src/train.py} (sans l'option \texttt{--skip-mlflow}), le script enregistre:
\begin{itemize}
  \item \textbf{Param\`etres}: nom du mod\`ele, strat\'egie (Class\_Weight ou SMOTE), indicateur \texttt{use\_smote}.
  \item \textbf{M\'etriques}: accuracy, precision, recall, f1\_score, roc\_auc.
  \item \textbf{Mod\`ele entra\^in\'e}: sauvegarde du mod\`ele en artifact (scikit-learn / XGBoost / LightGBM).
\end{itemize}

\subsection{O\`u sont stock\'ees les exp\'eriences}
Les runs MLflow sont stock\'es en local dans le dossier:\newline
\texttt{mlruns/}

\subsection{Lancer l'interface MLflow}
Pour visualiser et comparer les exp\'eriences:
\begin{verbatim}
mlflow ui
\end{verbatim}
Puis ouvrir l'URL affich\'ee (souvent \url{http://127.0.0.1:5000}).

\subsection{Compl\'ementarit\'e avec les sorties sur disque}
En plus de MLflow, ce projet sauvegarde syst\'ematiquement les artifacts dans \texttt{outputs/runs/} (m\^eme si MLflow est d\'esactiv\'e). Cela permet:
\begin{itemize}
  \item de consulter facilement les m\'etriques sous forme CSV/JSON,
  \item de r\'ecup\'erer les matrices de confusion et rapports de classification,
  \item de d\'eployer un mod\`ele sp\'ecifique via son \texttt{model.pkl}.
\end{itemize}

\section{Exp\'eriences et r\'esultats}
\subsection{Sauvegarde des sorties}
Chaque run est sauvegard\'e dans:
\begin{verbatim}
outputs/runs/<RUN_NAME>/<TIMESTAMP>/
\end{verbatim}
avec:
\begin{itemize}
  \item \texttt{metrics.json}, \texttt{metrics.csv}
  \item \texttt{confusion\_matrix.csv}
  \item \texttt{classification\_report.csv}
  \item \texttt{model.pkl}
\end{itemize}

Un r\'esum\'e global est disponible dans \texttt{outputs/runs/runs\_summary.csv}.

\subsection{Tableau comparatif des mod\`eles (run du 14/02/2026)}
Les valeurs ci-dessous ont \`et\'e extraites de \texttt{outputs/runs/runs\_summary.csv}.

\begin{table}[H]
\centering
\small
\begin{tabular}{lccccc}
\toprule
\textbf{Run} & \textbf{Acc.} & \textbf{Prec.} & \textbf{Rec.} & \textbf{F1} & \textbf{ROC-AUC} \\
\midrule
Logistic\_Regression\_Class\_Weight & 0.7388 & 0.5051 & 0.7968 & 0.6183 & 0.8398 \\
Logistic\_Regression\_SMOTE & 0.7424 & 0.5095 & 0.7888 & 0.6191 & 0.8391 \\
Decision\_Tree\_Class\_Weight & 0.7189 & 0.4813 & 0.7567 & 0.5884 & 0.7780 \\
Decision\_Tree\_SMOTE & 0.7410 & 0.5090 & 0.6791 & 0.5819 & 0.7913 \\
Random\_Forest\_Class\_Weight & 0.7764 & 0.5675 & 0.6631 & 0.6116 & 0.8348 \\
Random\_Forest\_SMOTE & 0.7679 & 0.5528 & 0.6578 & 0.6007 & 0.8325 \\
XGBoost\_Class\_Weight & 0.7480 & 0.5178 & 0.7380 & 0.6086 & 0.8314 \\
XGBoost\_SMOTE & 0.7289 & 0.4932 & 0.7754 & 0.6029 & 0.8321 \\
LightGBM\_Class\_Weight & 0.7566 & 0.5283 & 0.7727 & 0.6276 & 0.8363 \\
LightGBM\_SMOTE & 0.7793 & 0.5763 & 0.6364 & 0.6048 & 0.8365 \\
SVM\_Class\_Weight & 0.7417 & 0.5088 & 0.7701 & 0.6128 & 0.8161 \\
SVM\_SMOTE & 0.7388 & 0.5058 & 0.7005 & 0.5874 & 0.8088 \\
\bottomrule
\end{tabular}
\caption{Comparaison des runs (14/02/2026).}
\end{table}

\subsection{Meilleur mod\`ele}
Le meilleur mod\`ele selon le \textbf{F1-score} est:
\begin{itemize}
  \item \textbf{LightGBM\_Class\_Weight}
  \item Accuracy: 0.7566
  \item Precision: 0.5283
  \item Recall: 0.7727
  \item F1-score: 0.6276
  \item ROC-AUC: 0.8363
\end{itemize}

\section{D\'eploiement et inf\'erence (Flask)}
\subsection{Objectif}
L'objectif du d\'eploiement est de permettre \`a un utilisateur (ou un agent m\'etier) de saisir les informations d'un client via une interface web, puis d'obtenir une pr\'ediction instantan\'ee.

\subsection{Fonctionnement}
Le serveur Flask charge:
\begin{itemize}
  \item Le meilleur mod\`ele: \texttt{models/best\_model.pkl}
  \item Les artifacts de preprocessing sauvegard\'es: \texttt{outputs/preprocessed/scaler.pkl} et \texttt{outputs/preprocessed/label\_encoders.pkl}
\end{itemize}

Deux modes sont disponibles:
\begin{itemize}
  \item \textbf{Interface web}: page \texttt{/} (formulaire) + soumission \texttt{/predict-form}.
  \item \textbf{API JSON}: endpoint \texttt{POST /predict} pour int\'egration.
\end{itemize}

\subsection{Ex\'ecution}
\textbf{Local}:
\begin{verbatim}
python src/train.py --only-preprocess --save-preprocessed
python src/train.py --skip-mlflow
python app.py
\end{verbatim}

\textbf{Production} (WSGI Waitress):
\begin{verbatim}
python -m waitress --host=0.0.0.0 --port=5000 app:app
\end{verbatim}

\section{Discussion et recommandations}
\subsection{Analyse rapide}
Les mod\`eles de boosting (LightGBM, XGBoost) obtiennent des performances globalement \`a la hauteur des attentes. Dans ce run, LightGBM avec pond\'eration des classes maximise le F1-score.

\subsection{Recommandations business}
\begin{itemize}
  \item Cibler les clients \`a risque (probabilit\'e \'elev\'ee) avec des offres de r\'etention.
  \item Porter une attention particuli\`ere aux contrats \emph{month-to-month}.
  \item Am\'eliorer le support technique et la proposition de valeur pour les services internet.
\end{itemize}

\section{Conclusion}
Ce projet a permis de construire un pipeline complet de pr\'ediction du churn, depuis le pr\'etraitement des donn\'ees jusqu'au d\'eploiement d'un outil d'inf\'erence web. La comparaison multi-mod\`eles a montr\'e que \textbf{LightGBM (Class\_Weight)} est le meilleur mod\`ele selon le F1-score pour l'ex\'ecution du 14/02/2026. Les artifacts et m\'etriques sauvegard\'es facilitent la reproductibilit\'e, l'analyse et l'industrialisation.

\end{document}
